\section{Derived numerical cutoffs for ordinary alternating tensor factors}
\label{sec:effective-factorwise-tail-cutoffs}

\paragraph{Status and exact scope.}
HD-FT16--17, reserved question Q2dq, is a derived quantitative
corollary by \texttt{/root/g34\_eightfolds}.  It introduces no new
character or genus theorem.  The initial conditional proof and its first
exact certificate are retained unchanged.  HD-GC06, equation (EG11),
was subsequently independently reviewed and accepted at 02:38:42 UTC
on 22 November 2026; its accepted numerical conclusion is now used.
Complete author and separate root whole-proof readings pass, together with the complete source/status, checker and first-output readings. The first conditional certificate and subsequent status-only update are both retained; no duplicate arithmetic execution is claimed.  All character estimates and
actual quotient-cover interfaces used here are the already reviewed
HD-FT01--12.  No accepted file is changed by this addendum.

Throughout retain the actual factorwise setup (FT.1)--(FT.2): a
connected finite cover of the sphere, an actual quotient onto \(A_n\)
or \(S_n\), an actual factorization \(V=W\otimes U\) up to scalars,
and the genuine generating product-one inertia budget
\[
 e=\dim V\ge a=\dim W,\qquad
                \sum_i\operatorname{rank}(1-T_i)=2e+6.
 \tag{FTE.1}
\]
There is no lower bound on the number of genuine projected marks.
The tensor partner, its degree, its field and every scalar are arbitrary.
The genus premise being substituted is precisely
\[
 g(X_3)>n^2/288\qquad(n\ge10^6),
 \tag{FTE.2}
\]
where \(X_3\) is the actual connected quotient on unordered triples.
This is the stated conclusion of HD-GC06, not a proposed numerical
interpretation of an unspecified published asymptotic constant.

\subsection*{HD-FT16. Effective small-tail cutoffs with every error retained}

For a partition \(\mu\vdash k\), write the reviewed character
polynomial as \(q_\mu=C f^k+R\), with \(C>0\), and put
\[
 M=1+\|R\|_1/C,\qquad
 G=\left\lceil2M(M+1)+2\right\rceil.
 \tag{FTE.3}
\]
The norm is the sum of absolute monomial coefficients.  The integer
values below are recomputed from the retained checked polynomials;
they are not estimates of genus constants.

\begin{proposition}[HD-FT16]
Assume (FTE.2).  The following ordinary \(S_n\)-labels, their
transposes and their irreducible \(A_n\)-restrictions cannot occur
in (FTE.1):
\[
 \begin{array}{c|c}
 k&\text{sufficient integer natural-degree threshold}\\ \hline
 3& n\ge530841600,\\
 4& n\ge\max\{10^6,\lceil256M\rceil,\lceil1920M^2\rceil\},\\
 5,6&n\ge\max\{10^6,\lceil256M\rceil,
                              \lceil\sqrt{288G}\rceil\}.
 \end{array}
 \tag{FTE.4}
\]
In particular the single threshold \(530841600\) excludes every
tail of size three through six in this actual factorwise setup.
\end{proposition}

\begin{proof}
For tails five and six, HD-FT07--08 supplies an error-free cubic
absolute-character bound in the high-fixed-point region and error
\(2M(M+1)/n^3\) on at most ten remaining entries.  Its proof uses
\(n\ge256M\), \(a\ge2n^3\), and only the condition
\(g(X_3)>G\) in its geometric step.  Equation (FTE.2) supplies that
condition once \(n^2\ge288G\), including equality because the genus
bound is strict.  This gives the last row of (FTE.4), with the entire
\(6/e\) term already retained in HD-FT08.

For tail four, use the local estimates of HD-FT09: for \(n\ge256M\)
one has \(a\ge2n^3\), and at most ten entries have absolute-character
error \(4M^2/n^2\) compared with the natural cube.  The ordered
distinct-triple quotient maps to \(X_3\) with degree six.  Thus its
genus \(\gamma_3\) satisfies
\(\gamma_3-1\ge6(g(X_3)-1)\), whence
\(\gamma_3>n^2/48-5\).  The actual five-orbit decomposition of the
full cube gives
\[
 \sum_i\operatorname{ind}(x_i,[n]^3)
   >2n^3+n^2/24-20.
 \tag{FTE.5}
\]
After subtracting the ten errors and dividing by \(n^3\),
\[
 \frac1e\sum_i\operatorname{rank}(1-T_i)
 >2+\frac1{24n}-\frac{40M^2}{n^2}-\frac{20}{n^3}
 \ge2+\frac1{48n}-\frac{20}{n^3}
 >2+\frac3{n^3}\ge2+\frac6e.
 \tag{FTE.6}
\]
The weak inequality uses \(n\ge1920M^2\).  The last strict one
uses \(n^2>1104\), which follows from \(n\ge10^6\).
This proves the middle row without any unevaluated \(c,d,n_0\).

For tail three, HD-FT11 gives \(a\ge n^3/12\) and an error
\(48/n^{3/2}\) relative to the unordered-triple permutation
character, on at most ten entries once \(n\ge1024\).
Put \(B=\binom n3\).  Exact Riemann--Hurwitz on \(X_3\), together
with (FTE.2) and \(B\le n^3/6\), yields
\[
 \begin{split}
 \frac1e\sum_i\operatorname{rank}(1-T_i)
 &>2+\frac1{24n}-\frac{12}{n^3}-\frac{480}{n^{3/2}}\\
 &\ge2+\frac1{48n}-\frac{12}{n^3}
 >2+\frac{72}{n^3}\ge2+\frac6e.
 \end{split}
 \tag{FTE.7}
\]
For the first line keep the positive numerator together:
\((2g(X_3)-2)/B>(n^2/144-2)/B
\ge1/(24n)-12/n^3\), because \(n\ge10^6\) and
\(B\le n^3/6\).  No separate upper bound on \(2/B\) is used.
The second line uses
\(n\ge23040^2=530841600\), so \(480/\sqrt n\le1/48\).
Its strict inequality follows from \(n^2>4032\).
This proves the first row.  All estimates concern every target
eigenvalue and survive arbitrary tensor spectral convolution,
exactly as in the cited local lemmas.
\end{proof}

\paragraph{Exact per-label arithmetic.}
For tail three the values of \(M\) are \(30,18,42\) for
\(\mu=(3),(2,1),(1,1,1)\), respectively.  The sharper direct
tail-three argument uses the common threshold in (FTE.4), rather
than these coefficient constants.  For the other small tails the
complete values and thresholds are
\[
\begin{array}{c|r|r}
\mu\vdash4&M&N_\mu\\\hline
(4)&192&70778880\\
(3,1)&96&17694720\\
(2,2)&120&27648000\\
(2,1,1)&104&20766720\\
(1,1,1,1)&288&159252480
\end{array}
\qquad
\begin{array}{c|r|r|r}
\mu\vdash5&M&G&N_\mu\\\hline
(5)&1320&3487442&1000000\\
(4,1)&570&650942&1000000\\
(3,2)&504&509042&1000000\\
(3,1,1)&460&424122&1000000\\
(2,2,1)&600&721202&1000000\\
(2,1,1,1)&690&953582&1000000\\
(1,1,1,1,1)&2280&10401362&1000000
\end{array}
 \tag{FTE.8}
\]
\[
\begin{array}{c|r|r|r}
\mu\vdash6&M&G&N_\mu\\\hline
(6)&11760&276618722&3010560\\
(5,1)&4464&39863522&1142784\\
(4,2)&3040&18489282&1000000\\
(4,1,1)&3312&21945314&1000000\\
(3,3)&4224&35692802&1081344\\
(3,2,1)&2340&10955882&1000000\\
(3,1,1,1)&3360&22585922&1000000\\
(2,2,2)&5040&50813282&1290240\\
(2,2,1,1)&4000&32008002&1024000\\
(2,1,1,1,1)&6048&73168706&1548288\\
(1,1,1,1,1,1)&21600&933163202&5529600
\end{array}
 \tag{FTE.9}
\]
The maximal \(M\) at tail sizes \(3,4,5,6\) is respectively
\(42,288,2280,21600\); the maximal threshold is respectively
\(530841600,159252480,1000000,5529600\).
All bounds are sufficient thresholds, not claims of optimality.

\subsection*{HD-FT17. A numerical all-ordinary-tail conclusion}

\begin{corollary}[HD-FT17]
Assume HD-GC06.  In the actual factorwise setup (FTE.1), any
ordinary symmetric or alternating irreducible factor with
\[
                         n\ge N_*:=241^{482}
 \tag{FTE.10}
\]
must be extendible from \(S_n\), and its partition label satisfies
\[
                \min\{n-\lambda_1,n-\lambda'_1\}\le2.
 \tag{FTE.11}
\]
The conclusion includes arbitrary factorwise tensor partners and
all coefficient fields.  It does not assert existence of the
remaining sizes zero, one or two, or a bounded partner catalogue.
\end{corollary}

\begin{proof}
HD-FT02 removes all extendible labels with both tails at least241,
and HD-FT03 removes every split alternating constituent when
\(n>3500000\).  HD-FT06 already removes tails seven through240 at
\(n\ge241^{482}\), with its explicit uniform coefficient comparison.
Only sizes three through six remain to be addressed.  They are
excluded by HD-FT16 because
\(241^{482}>200^4=1600000000>530841600\).
This proves the corollary.  The actual quotient/factorization
hypothesis is retained throughout; no abstract normal component is
silently substituted for it.
\end{proof}

\paragraph{Boundary of the numerical conclusion.}
This addendum removes the unspecified triple-genus constants from
the stated ordinary-tail exclusions using the accepted HD-GC06.
It does not evaluate the exact Monderer--Neftin pair-classification
constant \(N_{14}\), replace the seven-mark conclusion of HD-FT15
by an unproved effective version, bound arbitrary original overgroups,
or construct a source or a detecting cycle.  HD-GC07's separate
eight-branch consequence has its own weaker effective
four-partner-entry conclusion.  The two assertions must not be conflated.
The bound \(N_*\) is not substituted for FQ or TW in a core whose
quasisimple factors are permuted: the required stable-factor quotient
and tensor factorization would first have to be established there.
